% Generated draft: requires mathematical and source review.
\documentclass[11pt]{article}
\usepackage[margin=25mm]{geometry}
\usepackage{fontspec}
\setmainfont{Latin Modern Roman}
\usepackage{amsmath,amssymb,mathtools,unicode-math}
\setmathfont{Latin Modern Math}
\usepackage{enumitem,needspace,xurl,hyperref}
\setlength{\parindent}{0pt}
\setlength{\parskip}{5pt}
\setlength{\emergencystretch}{3em}
\newcommand{\N}{\mathbb{N}}
\newcommand{\Z}{\mathbb{Z}}
\newcommand{\Q}{\mathbb{Q}}
\newcommand{\R}{\mathbb{R}}
\newcommand{\C}{\mathbb{C}}
\newcommand{\F}{\mathbb{F}}
\newcommand{\A}{\mathbb{A}}
\newcommand{\PP}{\mathbb{P}}
\newcommand{\E}{\mathbb{E}}
\newcommand{\eps}{\varepsilon}
\newcommand{\dd}{\,\mathrm{d}}
\begin{document}
\section*{UnsolvedMath ID: 30000346}
\textbf{Source problem code:} catalogue entry
\subsection{Definitions and mathematical statement}
\textbf{Two-part selection problem.} For finite-volume corrections to the Lifshitz--Slyozov--Wagner (LSW) coarsening model, determine whether the self-consistent theory closed at the two-particle distribution level resolves the weak selection problem for self-similar asymptotic states. Determine also the role of particle collisions in selection. Both the validity of the proposed closure and the collision contribution to long-time selection are part of the request.\par The LSW picture describes coarsening through a distribution of particle sizes. A one-particle distribution records sizes and, when needed, positions; a two-particle distribution additionally records pair correlations and relative configuration. Closing at that level means specifying a rule for the higher-order quantities needed in the evolution equations in terms of the retained distributions. The exact self-consistent equations and closure rule are not supplied, so factorization, independence, or a particular collision kernel is not imposed silently. The phrase \emph{finite-volume corrections} may refer to a nonzero particle volume fraction rather than a bounded physical container; the intended meaning is not defined and is not identified with a numerical finite-volume scheme.\par A self-similar asymptotic state means a limiting profile after the model's prescribed rescaling of size, density, and possibly pair separation. In schematic notation a rescaled distribution \(\widetilde f_t\) should approach an admissible profile \(F_\theta\) in a specified topology. The weak selection question asks which members of the self-similar family arise from which initial data and whether the corrected model supplies the missing selection criterion. It does not automatically assert one universal profile independent of initial data. The profile family, scaling laws, data class, convergence topology, and order of the long-time and small-volume limits are absent. To determine collisions' role one must also specify what a collision does and how its rate enters the same model; switching off a term in an unspecified equation is not a defined comparison.\par The supplied triage checked 2026-08-21 labels the problem partially solved, reporting developed finite-volume collision corrections while describing selection and asymptotic closure as unresolved. No exact theorem, model, or closure validity result is reproduced. This dated assessment requires independent verification.
\subsection{Short English statement}
Does the specified pair-distribution closure select the self-similar states of a finite-volume LSW correction, and how do collisions affect that selection? The equations, closure, collision law, and limiting procedure are missing; selection need not mean a universal profile for every initial datum.
\subsection{Sources}
\begin{itemize}\item Local UnsolvedMath registry, permanent ID 30000346, code catalogue entry; \emph{Selection in Finite-Volume LSW Coarsening Models}, Oberwolfach Reports Open Problems.\item \emph{Reactive Flow and Transport Through Complex Systems (2006)}, Oberwolfach Reports, DOI 10.4171/owr/2005/49, pp. 2761--2832. \url{https://doi.org/10.4171/owr/2005/49}. Proposed year and year tag: 2005; no proposer named.\item Supplied literature label \emph{Oberwolfach Report 49}: \url{https://doi.org/10.4171/owr/2005/49}. No separate closure or collision theorem reference is supplied.\item Literature triage checked 2026-08-21, status \texttt{partially\_solved}; record updated 2026-09-16. Statement status \texttt{exact}, with confirmation reported without textual change. Verify dated claims independently.\end{itemize}
\end{document}
